\documentclass[11pt]{article}
\usepackage[a4paper,margin=22mm]{geometry}
\usepackage{fontspec}
\setmainfont{DejaVu Serif}
\setsansfont{DejaVu Sans}
\usepackage{microtype}
\usepackage{xcolor}
\usepackage{hyperref}
\usepackage{enumitem}
\usepackage{parskip}
\definecolor{Accent}{HTML}{971D32}
\definecolor{Muted}{HTML}{666666}
\hypersetup{colorlinks=true,urlcolor=Accent,linkcolor=Accent}
\setlist{leftmargin=1.6em,itemsep=0.3em,topsep=0.35em}
\setcounter{tocdepth}{2}
\renewcommand{\contentsname}{Contents}
\newcommand{\sectionrule}{\par\vspace{-0.5em}\noindent\textcolor{Accent}{\rule{\linewidth}{0.6pt}}\par}
\begin{document}

\begin{center}
{\sffamily\bfseries\color{Accent} TOEFL WORD OF THE DAY\par}
\vspace{8mm}
{\fontsize{42}{48}\selectfont\bfseries residual\par}
\vspace{2mm}
{\Large\sffamily\color{Accent} /rɪˈzɪdʒuəl/\par}
\vspace{2mm}
{\small\sffamily Local audio phrase: residual\par}
{\small\sffamily Audio file: audios/residual.mp3\par}
\vspace{2mm}
{\large\sffamily\color{Muted} adjective and countable noun\par}
\vspace{5mm}
{\sffamily remaining after something ends \textbullet\ an amount left over \textbullet\ difference from a predicted value\par}
\vspace{6mm}
{\large Residual describes or names something that remains after the main part has been removed, used, explained, or completed.\par}
\vspace{6mm}
\begin{minipage}{0.84\textwidth}
\itshape “The treatment reduced the infection, but some residual symptoms remained.”
\end{minipage}\par
\vspace{10mm}
\textbf{Date:} July 29th 2026\quad\textbullet\quad \textbf{Level:} B1--mid-B2 TOEFL
\vfill
Created and compiled by \textbf{Amir Kooshky}\par
\vspace{2mm}
\href{https://t.me/KooshkyTOEFL}{Telegram Channel}\quad\textbullet\quad
\href{https://www.instagram.com/kooshkytoefl}{Instagram}\quad\textbullet\quad
\href{https://t.me/KooshkyTOEFL\_pv}{Personal Telegram}
\end{center}
\newpage
\tableofcontents
\newpage

\section{Meanings and usage}
\sectionrule
\subsection{Meaning 1: Remaining after a process}
\textbf{Part of speech:} adjective

\textbf{IPA:} /rɪˈzɪdʒuəl/

\textbf{Pronunciation phrase:} residual

\textbf{Local audio file:} audios/residual.mp3

\textbf{Register:} formal; common in science, medicine, economics, engineering, and academic writing

\textbf{Definition:} remaining after most of something has disappeared, been removed, or ended

\textbf{In simpler words:} left over after most is gone

\textbf{Typical pattern:} residual + effect, amount, heat, moisture, chemical, symptom, or risk

\subsubsection{Natural examples}
\begin{enumerate}
  \item The treatment reduced the infection, but some residual symptoms remained.
  \item Residual heat from the machinery continued to warm the room.
  \item The soil contained residual chemicals from earlier farming practices.
  \item Researchers found residual traces of the substance in the water.
  \item The conflict left residual tension between the two communities.
  \item Residual moisture can damage the equipment if it is not removed.
  \item The economic crisis had residual effects that continued for several years.
\end{enumerate}

\subsubsection{Nuance and implication}
\textbf{Central nuance:} The word emphasizes that the main part has gone or been dealt with, but a smaller part still remains.

\textbf{Usual implication:} The remaining part may continue to have an effect even though it is relatively small.

\textbf{Compared with remaining:} Remaining is a general word for anything left. Residual is more formal and often suggests that something is left after a specific process, treatment, calculation, or event.

\subsubsection{Grammar notes}
\textbf{How to use it:} Use it before a noun to describe something that remains after a process or event.

\textbf{What can follow it:} It commonly comes before nouns such as effect, amount, heat, moisture, chemical, symptom, damage, risk, and value.

\textbf{Position in a sentence:} It is most often used before a noun, as in residual damage. Use remaining rather than residual in many ordinary everyday situations.

\subsubsection{Useful patterns}
\textbf{Pattern:} residual + effect or impact

\textbf{Use:} Use this for an effect that continues after its main cause or event has ended.

\textbf{Example:} The policy had residual effects on employment in the region.

\textbf{Pattern:} residual + substance or material

\textbf{Use:} Use this for a small amount of material left after cleaning, processing, or use.

\textbf{Example:} Residual oil was found inside the pipe.

\textbf{Pattern:} residual + symptom or damage

\textbf{Use:} Use this for a problem that remains after recovery or repair.

\textbf{Example:} Most patients recovered fully, although a few experienced residual pain.

\textbf{Pattern:} residual + value

\textbf{Use:} Use this for the estimated value that an asset still has after use or depreciation.

\textbf{Example:} The vehicle is expected to have a high residual value after three years.

\subsubsection{Common wrong usage}
\paragraph{Correction 1}
\textbf{Wrong:} There was a residual of tension after the meeting.

\textbf{Correct:} There was residual tension after the meeting.

\textbf{Why:} When residual describes a noun, place it directly before that noun.

\paragraph{Correction 2}
\textbf{Wrong:} The medicine caused some remained effects.

\textbf{Correct:} The medicine caused some residual effects.

\textbf{Why:} Use residual, not remained, as an adjective before a noun.

\paragraph{Correction 3}
\textbf{Wrong:} I ate the residual pizza for lunch.

\textbf{Correct:} I ate the leftover pizza for lunch.

\textbf{Why:} Residual is generally too formal and technical for ordinary food left after a meal. Leftover is more natural.

\subsection{Meaning 2: Amount left over}
\textbf{Part of speech:} countable noun

\textbf{IPA:} /rɪˈzɪdʒuəl/

\textbf{Pronunciation phrase:} residual

\textbf{Local audio file:} audios/residual.mp3

\textbf{Register:} formal and technical; used mainly in science, finance, mathematics, and manufacturing

\textbf{Definition:} an amount that remains after other amounts have been removed, used, or calculated

\textbf{In simpler words:} an amount left after the rest is removed

\textbf{Typical pattern:} a residual from + process or the residual after + calculation

\subsubsection{Natural examples}
\begin{enumerate}
  \item The researchers examined the residuals left after the chemical reaction.
  \item Any residual from the manufacturing process must be disposed of safely.
  \item The company receives the residual after all operating expenses have been paid.
  \item The filtering system removes most particles, but a small residual may remain.
  \item The accountant calculated the residual after subtracting the fixed costs.
  \item Scientists analyzed the residuals to identify materials that had not broken down.
  \item The residual was too small to affect the final result.
\end{enumerate}

\subsubsection{Nuance and implication}
\textbf{Central nuance:} As a noun, residual refers to the actual amount or quantity that is left, not merely its quality of remaining.

\textbf{Usual implication:} The amount is often determined after subtraction, processing, filtering, or distribution.

\subsubsection{Grammar notes}
\textbf{How to use it:} Use a residual for one remaining amount and residuals for several separate remaining amounts.

\textbf{What can follow it:} It may be followed by from to identify the process that produced it, or after to identify the calculation or deduction.

\textbf{Common preposition:} Say a residual from the process or the residual after expenses were paid.

\textbf{Noun use:} This is countable in technical use. However, residue is often more natural when referring to physical material left on a surface.

\subsubsection{Useful patterns}
\textbf{Pattern:} a residual from + process

\textbf{Use:} Use this for an amount left as a result of a process.

\textbf{Example:} The laboratory tested a residual from the purification process.

\textbf{Pattern:} the residual after + deduction or payment

\textbf{Use:} Use this for the amount remaining after costs or other amounts have been subtracted.

\textbf{Example:} Investors receive the residual after all debts have been paid.

\textbf{Pattern:} calculate the residual

\textbf{Use:} Use this when the remaining amount is found mathematically.

\textbf{Example:} The analyst calculated the residual by subtracting total costs from revenue.

\subsubsection{Common wrong usage}
\paragraph{Correction 1}
\textbf{Wrong:} The workers cleaned the residual from the table.

\textbf{Correct:} The workers cleaned the residue from the table.

\textbf{Why:} Residue is normally used for physical material left on a surface. Residual as a noun is more technical and often refers to a calculated or measured amount.

\paragraph{Correction 2}
\textbf{Wrong:} The company kept residual after paying its debts.

\textbf{Correct:} The company kept the residual after paying its debts.

\textbf{Why:} A singular countable use normally needs a, an, or the.

\subsection{Meaning 3: Difference from a prediction}
\textbf{Part of speech:} countable noun

\textbf{IPA:} /rɪˈzɪdʒuəl/

\textbf{Pronunciation phrase:} residual

\textbf{Local audio file:} audios/residual.mp3

\textbf{Register:} technical; common in statistics, data analysis, economics, and scientific research

\textbf{Definition:} in statistics, the difference between an observed value and the value predicted by a model

\textbf{In simpler words:} the gap between an actual and predicted value

\textbf{Typical pattern:} observed value minus predicted value equals residual

\subsubsection{Natural examples}
\begin{enumerate}
  \item A large residual may indicate that the model does not explain the observation well.
  \item The students calculated a residual for each data point.
  \item The residuals were plotted to check whether the statistical model was appropriate.
  \item A positive residual means that the observed value is higher than the predicted value.
  \item Several unusually large residuals suggested that the data contained outliers.
  \item The model was revised to reduce the size of the residuals.
  \item Randomly distributed residuals generally provide evidence that the model fits the data reasonably well.
\end{enumerate}

\subsubsection{Nuance and implication}
\textbf{Central nuance:} A residual measures how far one observation is from the value that a statistical model predicts.

\textbf{Usual implication:} Small and randomly distributed residuals often suggest a better-fitting model, while large or patterned residuals may reveal problems.

\subsubsection{Grammar notes}
\textbf{How to use it:} Use a residual for the difference associated with one observation and residuals for the differences across a data set.

\textbf{What can follow it:} It often appears with verbs such as calculate, examine, analyze, plot, minimize, and reduce.

\textbf{Noun use:} This is a countable technical noun. The plural residuals is especially common when discussing an entire data set.

\subsubsection{Useful patterns}
\textbf{Pattern:} calculate a residual

\textbf{Use:} Use this for finding the difference between one observed value and its predicted value.

\textbf{Example:} Calculate a residual for each point in the data set.

\textbf{Pattern:} analyze or plot the residuals

\textbf{Use:} Use this when examining whether a statistical model fits the data well.

\textbf{Example:} The researchers plotted the residuals before accepting the model.

\textbf{Pattern:} a positive or negative residual

\textbf{Use:} Use positive when the observation is above the prediction and negative when it is below.

\textbf{Example:} The final observation produced a large negative residual.

\subsubsection{Common wrong usage}
\paragraph{Correction 1}
\textbf{Wrong:} The residual is the predicted value itself.

\textbf{Correct:} The residual is the difference between the observed value and the predicted value.

\textbf{Why:} A residual is a difference, not the predicted value.

\paragraph{Correction 2}
\textbf{Wrong:} The researcher calculated the residuals between the two groups.

\textbf{Correct:} The researcher calculated the residuals for the observations.

\textbf{Why:} Statistical residuals are normally calculated for observations by comparing observed and predicted values.

\section{Word family}
\sectionrule
\subsection{residue}
\textbf{Part of speech:} countable or uncountable noun

\textbf{IPA:} /ˈrɛzəˌduː/

\textbf{Definition:} a substance or small amount of material that remains after a process such as burning, evaporation, or cleaning

\textbf{Relationship to the headword:} Residue names material that remains, while residual commonly describes something that remains.

\textbf{Grammar pattern:} residue on + surface or residue from + process

\textbf{Usage note:} Use residue without a number for material in general and residues for separate types or deposits.

\subsubsection{Examples}
\begin{enumerate}
  \item The cleaner left a sticky residue on the glass.
  \item Scientists tested the soil for pesticide residues.
\end{enumerate}

\subsection{residually}
\textbf{Part of speech:} adverb

\textbf{IPA:} /rɪˈzɪdʒuəli/

\textbf{Definition:} in a way that remains after the main part has been removed, completed, or explained

\textbf{Relationship to the headword:} This adverb describes something occurring or existing as a residual.

\textbf{Usage note:} It is uncommon and mainly appears in specialized academic or technical writing.

\subsubsection{Examples}
\begin{enumerate}
  \item The substance was residually present in the soil.
  \item The remaining funds were residually allocated among the partners.
\end{enumerate}

\section{Common collocations}
\sectionrule
\subsection{residual effect}
\textbf{Applies to:} Remaining after a process

\textbf{Meaning:} an effect that continues after the main action or event has ended

\textbf{Pattern:} residual effect of + event, treatment, or substance

\textbf{Example:} The medication may have residual effects for several hours.

\subsection{residual impact}
\textbf{Applies to:} Remaining after a process

\textbf{Meaning:} an influence that continues after the main event

\textbf{Example:} The recession had a residual impact on household spending.

\subsection{residual risk}
\textbf{Applies to:} Remaining after a process

\textbf{Meaning:} risk that remains after efforts have been made to reduce it

\textbf{Example:} Even after the repairs, a small residual risk remained.

\subsection{residual heat}
\textbf{Applies to:} Remaining after a process

\textbf{Meaning:} heat that remains after a machine or process has stopped

\textbf{Example:} The system uses residual heat to warm nearby buildings.

\subsection{residual moisture}
\textbf{Applies to:} Remaining after a process

\textbf{Meaning:} a small amount of water that remains after drying

\textbf{Example:} Residual moisture inside the container encouraged mold growth.

\subsection{residual value}
\textbf{Applies to:} Remaining after a process

\textbf{Meaning:} the value an asset is expected to retain after use or depreciation

\textbf{Pattern:} the residual value of + asset

\textbf{Example:} The machine's residual value was estimated at ten thousand dollars.

\subsection{residual income}
\textbf{Applies to:} Amount left over

\textbf{Meaning:} income remaining after specified costs or required returns are subtracted

\textbf{Example:} Residual income is calculated after required expenses and returns have been deducted.

\subsection{calculate the residual}
\textbf{Applies to:} Amount left over

\textbf{Meaning:} determine the amount that remains after subtraction

\textbf{Example:} The accountant calculated the residual after deducting all liabilities.

\subsection{large residual}
\textbf{Applies to:} Difference from a prediction

\textbf{Meaning:} a large gap between an observed and predicted value

\textbf{Example:} A large residual showed that the prediction was far from the observed value.

\subsection{residual analysis}
\textbf{Applies to:} Difference from a prediction

\textbf{Meaning:} the examination of residuals to judge how well a statistical model fits

\textbf{Example:} The researchers used residual analysis to evaluate the model.

\subsection{plot the residuals}
\textbf{Applies to:} Difference from a prediction

\textbf{Meaning:} display residual values on a graph

\textbf{Pattern:} plot the residuals against + variable

\textbf{Example:} The students plotted the residuals to identify patterns in the model's errors.

\section{Synonyms and antonyms}
\sectionrule
\subsection{Close synonyms}
\subsubsection{remaining}
\textbf{Part of speech:} adjective

\textbf{Applies to:} Remaining after a process

\textbf{Shared meaning:} Both words describe something that is still left.

\textbf{Difference:} Remaining is more general and common in everyday language. Residual is more formal and often refers to something left after a particular process or event.

\textbf{Example:} The remaining funds will be used for maintenance.

\subsubsection{leftover}
\textbf{Part of speech:} adjective

\textbf{Applies to:} Remaining after a process

\textbf{Shared meaning:} Both can describe something that remains after the main part has been used.

\textbf{Difference:} Leftover is more informal and is especially common for food, materials, or unused items. Residual is more technical and academic.

\textbf{Example:} The leftover materials were stored for future use.

\subsubsection{remainder}
\textbf{Part of speech:} countable noun

\textbf{Applies to:} Amount left over

\textbf{Shared meaning:} Both words can refer to an amount left after subtraction or use.

\textbf{Difference:} Remainder is a common general word for what is left. Residual is more formal and often refers to an amount produced by a calculation or technical process.

\textbf{Example:} The remainder of the budget was transferred to the next year.

\subsubsection{error}
\textbf{Part of speech:} countable noun

\textbf{Applies to:} Difference from a prediction

\textbf{Shared meaning:} Both can describe a difference between a value and a reference point.

\textbf{Difference:} In statistics, error may refer to the difference between an observed value and the true underlying value, which may be unknown. A residual is the observable difference between an observed value and a model's predicted value.

\textbf{Example:} Measurement error can reduce the accuracy of an experiment.

\section{Words learners may confuse}
\sectionrule
\subsection{residue}
\textbf{Part of speech:} adjective versus noun

\textbf{Why learners confuse them:} The words share the same root and both involve the idea of something being left behind.

\textbf{Key difference:} Residual is usually an adjective meaning remaining, while residue is usually a noun referring to physical material left after a process.

\textbf{residual:} The laboratory detected residual chemicals in the water.

\textbf{residue:} A white residue remained at the bottom of the container.

\section{Etymology}
\sectionrule
\textbf{Origin:} Residual ultimately comes from Latin residuum, meaning something left behind or remaining.

\section{Practice exercises}
\sectionrule
\subsection{Word family choice}
Choose the best member of the word family for each blank.

\begin{enumerate}
  \item The cleaning product left a thin \_\_\_\_\_ on the surface.
  \item The patient experienced some \_\_\_\_\_ discomfort after the treatment.
\end{enumerate}

\subsection{Correct or incorrect?}
Decide whether each sentence is correct. Correct the inaccurate ones.

\begin{enumerate}
  \item The fire produced enough residual heat to warm the building.
  \item A sticky residual covered the surface.
  \item The residual is the difference between the observed and predicted values.
  \item The company hopes to eliminate every residual risk completely.
\end{enumerate}

\subsection{Collocation practice}
Complete each sentence with a natural collocation.

\begin{enumerate}
  \item The medicine may cause \_\_\_\_\_ effects even after treatment ends.
  \item The engineer measured the \_\_\_\_\_ moisture inside the storage tank.
  \item The analyst plotted the \_\_\_\_\_ to evaluate the statistical model.
\end{enumerate}

\subsection{Complete answer key}
\subsubsection{Word family choice}
\begin{enumerate}
  \item \textbf{residue.} The cleaning product left a thin residue on the surface. \emph{Use the noun residue for physical material left on a surface.}
  \item \textbf{residual.} The patient experienced some residual discomfort after the treatment. \emph{Use the adjective residual before the noun discomfort.}
\end{enumerate}

\subsubsection{Correct or incorrect?}
\begin{enumerate}
  \item \textbf{Correct} \emph{Residual heat is an established expression for heat that remains after a process.}
  \item \textbf{Incorrect}. A sticky residue covered the surface. \emph{Use residue for physical material left on a surface.}
  \item \textbf{Correct} \emph{This is the standard statistical meaning of residual.}
  \item \textbf{Correct} \emph{Residual risk naturally refers to risk that remains after preventive measures.}
\end{enumerate}

\subsubsection{Collocation practice}
\begin{enumerate}
  \item \textbf{residual.} The medicine may cause residual effects even after treatment ends. \emph{Residual effects are effects that continue after the main action has ended.}
  \item \textbf{residual.} The engineer measured the residual moisture inside the storage tank. \emph{Residual moisture is water that remains after drying or removal.}
  \item \textbf{residuals.} The analyst plotted the residuals to evaluate the statistical model. \emph{Plot the residuals is a standard expression in statistical analysis.}
\end{enumerate}

\vfill
\begin{center}
\small\color{Muted} Created and compiled by \textbf{Amir Kooshky}\\
\href{https://t.me/KooshkyTOEFL}{Telegram Channel}\quad\textbullet\quad
\href{https://www.instagram.com/kooshkytoefl}{Instagram}\quad\textbullet\quad
\href{https://t.me/KooshkyTOEFL\_pv}{Personal Telegram}
\end{center}
\end{document}
